% GENERATED FILE. Edit canonical lessons, then run npm run generate:books.
% canonical-source-hash: fnv1a64:ddb534121c8151db
% canonical-lessons: ML-C443-badhikkuka, ML-C443-svadhinikkuka, ML-C443-badal, ML-C443-paridhi, ML-C443-samayaparidhi

\chapter{Effects and Limits}
\label{ch:ml-effects-and-limits}
\hlchaptermodality{443}

\begin{chapteropening}
\textbf{By the end of this chapter:} \emph{I can say to affect, to influence, an alternative, a limit and a deadline.}

Complete the last lesson of chapter 443: I can say to affect, to influence, an alternative, a limit and a deadline.
\end{chapteropening}

\section[\texorpdfstring{\ml{ബാധിക്കുക} (bādhikkuka)}{ബാധിക്കുക (bādhikkuka)}]{\ml{ബാധിക്കുക} (bādhikkuka) — to affect}

\label{lesson:ML-C443-badhikkuka}



\begin{quote}
Before the new one: say the Malayalam for to estimate, then the Malayalam for a consequence.
\end{quote}

\subsection*{You'll want to know: \ml{ബാധിക്കുക}}
\textbf{\ml{ബാധിക്കുക}} — \emph{bādhikkuka} — \textquotedblleft{}to affect\textquotedblright{}.

\begin{quote}
\ml{പനി} \ml{ഗ്രാമത്തിലെ} \ml{പല} \ml{കുട്ടികളെയും} \ml{ബാധിച്ചു}.
\end{quote}

\begin{quote}
\emph{pani grāmattile pala kuṭṭikaḷeyuṁ bādhiccu.}
\end{quote}

\begin{quote}
The fever affected many children in the village.
\end{quote}

Past: \ml{ബാധിച്ചു}, \emph{bādhiccu}.

What is affected takes -\ml{നെ}. Usually something bad: an illness, a flood, a price rise.

\begin{cousinweb}
\textbf{\ml{ബാധിക്കുക}} (\emph{bādhikkuka}) is built from the Sanskrit root \emph{bādh}, \textquotedblleft{}to press hard, to harass\textquotedblright{}, whose noun \emph{bādha} is \textquotedblleft{}affliction, hurt, damage\textquotedblright{} (Monier-Williams), with the Malayalam verb ending -\ml{ിക്കുക}.
\end{cousinweb}

\begin{grammarlens}[title={What you've built}]
One more word for saying what happens if.
\end{grammarlens}

\subsection*{Your turn}
\begin{itemize}
\raggedright
  \item \emph{Say it:} \emph{bādhikkuka}
  \item \emph{Say it:} \emph{bādhikkuka}, once more
  \item \emph{Say it:} \emph{pani grāmattile pala kuṭṭikaḷeyuṁ bādhiccu.} — The fever affected many children in the village.
  \item \emph{Recall:} say the Malayalam for to estimate, then the Malayalam for a consequence, then say \emph{bādhikkuka} again
\end{itemize}

\begin{tcolorbox}[breakable,colback=teal!4,colframe=teal!35!black,title={Before you move on}]
What does \ml{ബാധിക്കുക} mean? (\textquotedblleft{}To affect\textquotedblright{}.) Say it once more.
\end{tcolorbox}
\section[\texorpdfstring{\ml{സ്വാധീനിക്കുക} (svādhīnikkuka)}{സ്വാധീനിക്കുക (svādhīnikkuka)}]{\ml{സ്വാധീനിക്കുക} (svādhīnikkuka) — to influence}

\label{lesson:ML-C443-svadhinikkuka}



\begin{quote}
Before the new one: say the Malayalam for a consequence, then the Malayalam for to affect.
\end{quote}

\subsection*{You'll want to know: \ml{സ്വാധീനിക്കുക}}
\textbf{\ml{സ്വാധീനിക്കുക}} — \emph{svādhīnikkuka} — \textquotedblleft{}to influence\textquotedblright{}.

\begin{quote}
\ml{അമ്മയുടെ} \ml{വാക്കുകൾ} \ml{അവന്റെ} \ml{തീരുമാനത്തെ} \ml{സ്വാധീനിച്ചു}.
\end{quote}

\begin{quote}
\emph{ammayuṭe vākkukaḷ avanṟe tīrumānatte svādhīniccu.}
\end{quote}

\begin{quote}
Mother's words influenced his decision.
\end{quote}

Past: \ml{സ്വാധീനിച്ചു}, \emph{svādhīniccu}.

The noun is \ml{സ്വാധീനം}, influence.

\begin{cousinweb}
\textbf{\ml{സ്വാധീനിക്കുക}} (\emph{svādhīnikkuka}) is built from Sanskrit \emph{svādhīna}, \textquotedblleft{}dependent on oneself, under one's own control\textquotedblright{} (Monier-Williams), with the Malayalam verb ending -\ml{ിക്കുക}: to bring a thing under one's sway.
\end{cousinweb}

\begin{grammarlens}[title={What you've built}]
One more word for saying what happens if.
\end{grammarlens}

\subsection*{Your turn}
\begin{itemize}
\raggedright
  \item \emph{Say it:} \emph{svādhīnikkuka}
  \item \emph{Say it:} \emph{svādhīnikkuka}, once more
  \item \emph{Say it:} \emph{ammayuṭe vākkukaḷ avanṟe tīrumānatte svādhīniccu.} — Mother's words influenced his decision.
  \item \emph{Recall:} say the Malayalam for a consequence, then the Malayalam for to affect, then say \emph{svādhīnikkuka} again
\end{itemize}

\begin{tcolorbox}[breakable,colback=teal!4,colframe=teal!35!black,title={Before you move on}]
What does \ml{സ്വാധീനിക്കുക} mean? (\textquotedblleft{}To influence\textquotedblright{}.) Say it once more.
\end{tcolorbox}
\section[\texorpdfstring{\ml{ബദൽ} (badal)}{ബദൽ (badal)}]{\ml{ബദൽ} (badal) — an alternative; a substitute}

\label{lesson:ML-C443-badal}



\begin{quote}
Before the new one: say the Malayalam for to affect, then the Malayalam for to influence.
\end{quote}

\subsection*{You'll want to know: \ml{ബദൽ}}
\textbf{\ml{ബദൽ}} — \emph{badal} — \textquotedblleft{}an alternative; a substitute\textquotedblright{}.

\begin{quote}
\ml{ആ} \ml{പദ്ധതിക്ക്} \ml{ഒരു} \ml{ബദൽ} \ml{കണ്ടെത്തണം}.
\end{quote}

\begin{quote}
\emph{ā paddhatikkŭ oru badal kaṇṭettaṇaṁ.}
\end{quote}

\begin{quote}
An alternative to that plan must be found.
\end{quote}

As an adjective it comes first: \ml{ബദൽ} \ml{വഴി}, an alternative route.

\begin{cousinweb}
\textbf{\ml{ബദൽ}} (\emph{badal}) is usually traced to Hindustani \emph{badal}, from Arabic \emph{badal}, \textquotedblleft{}a substitute, an exchange\textquotedblright{}.
\end{cousinweb}

\begin{grammarlens}[title={What you've built}]
One more word for saying what happens if.
\end{grammarlens}

\subsection*{Your turn}
\begin{itemize}
\raggedright
  \item \emph{Say it:} \emph{badal}
  \item \emph{Say it:} \emph{badal}, once more
  \item \emph{Say it:} \emph{ā paddhatikkŭ oru badal kaṇṭettaṇaṁ.} — An alternative to that plan must be found.
  \item \emph{Recall:} say the Malayalam for to affect, then the Malayalam for to influence, then say \emph{badal} again
\end{itemize}

\begin{tcolorbox}[breakable,colback=teal!4,colframe=teal!35!black,title={Before you move on}]
What does \ml{ബദൽ} mean? (\textquotedblleft{}An alternative; a substitute\textquotedblright{}.) Say it once more.
\end{tcolorbox}
\section[\texorpdfstring{\ml{പരിധി} (paridhi)}{പരിധി (paridhi)}]{\ml{പരിധി} (paridhi) — a limit}

\label{lesson:ML-C443-paridhi}



\begin{quote}
Before the new one: say the Malayalam for to influence, then the Malayalam for an alternative.
\end{quote}

\subsection*{You'll want to know: \ml{പരിധി}}
\textbf{\ml{പരിധി}} — \emph{paridhi} — \textquotedblleft{}a limit\textquotedblright{}.

\begin{quote}
\ml{എല്ലാത്തിനും} \ml{ഒരു} \ml{പരിധി} \ml{ഉണ്ട്}.
\end{quote}

\begin{quote}
\emph{ellāttinuṁ oru paridhi uṇṭŭ.}
\end{quote}

\begin{quote}
Everything has a limit.
\end{quote}

Within the limit: \ml{പരിധിക്കുള്ളിൽ}.

\begin{cousinweb}
\textbf{\ml{പരിധി}} (\emph{paridhi}) is borrowed from Sanskrit \emph{paridhi}, \textquotedblleft{}an enclosure, a fence, a wall\textquotedblright{} (Monier-Williams): a limit is the fence round a thing.
\end{cousinweb}

\begin{grammarlens}[title={What you've built}]
One more word for saying what happens if.
\end{grammarlens}

\subsection*{Your turn}
\begin{itemize}
\raggedright
  \item \emph{Say it:} \emph{paridhi}
  \item \emph{Say it:} \emph{paridhi}, once more
  \item \emph{Say it:} \emph{ellāttinuṁ oru paridhi uṇṭŭ.} — Everything has a limit.
  \item \emph{Recall:} say the Malayalam for to influence, then the Malayalam for an alternative, then say \emph{paridhi} again
\end{itemize}

\begin{tcolorbox}[breakable,colback=teal!4,colframe=teal!35!black,title={Before you move on}]
What does \ml{പരിധി} mean? (\textquotedblleft{}A limit\textquotedblright{}.) Say it once more.
\end{tcolorbox}
\section[\texorpdfstring{\ml{സമയപരിധി} (samayaparidhi)}{സമയപരിധി (samayaparidhi)}]{\ml{സമയപരിധി} (samayaparidhi) — a deadline; a time limit}

\label{lesson:ML-C443-samayaparidhi}



\begin{quote}
Before the new one: say the Malayalam for an alternative, then the Malayalam for a limit.
\end{quote}

\subsection*{You'll want to know: \ml{സമയപരിധി}}
\textbf{\ml{സമയപരിധി}} — \emph{samayaparidhi} — \textquotedblleft{}a deadline; a time limit\textquotedblright{}.

\begin{quote}
\ml{അപേക്ഷ} \ml{നൽകാനുള്ള} \ml{സമയപരിധി} \ml{നാളെ} \ml{തീരും}.
\end{quote}

\begin{quote}
\emph{apēkṣa nalkānuḷḷa samayaparidhi nāḷe tīruṁ.}
\end{quote}

\begin{quote}
The deadline for applying ends tomorrow.
\end{quote}

\ml{സമയം}, time, and \ml{പരിധി}, a limit.

\begin{cousinweb}
\textbf{\ml{സമയപരിധി}} (\emph{samayaparidhi}) is a Malayalam compound built from Sanskrit \emph{samaya}, \textquotedblleft{}a meeting; an appointed time\textquotedblright{} (CDIAL 13185), and \emph{paridhi}, \textquotedblleft{}an enclosure\textquotedblright{} (Monier-Williams): the fence round a stretch of time.
\end{cousinweb}

\begin{grammarlens}[title={What you've built}]
One more word for saying what happens if.
\end{grammarlens}

\subsection*{Your turn}
\begin{itemize}
\raggedright
  \item \emph{Say it:} \emph{samayaparidhi}
  \item \emph{Say it:} \emph{samayaparidhi}, once more
  \item \emph{Say it:} \emph{apēkṣa nalkānuḷḷa samayaparidhi nāḷe tīruṁ.} — The deadline for applying ends tomorrow.
  \item \emph{Recall:} say the Malayalam for an alternative, then the Malayalam for a limit, then say \emph{samayaparidhi} again
\end{itemize}

\begin{tcolorbox}[breakable,colback=teal!4,colframe=teal!35!black,title={Before you move on}]
What does \ml{സമയപരിധി} mean? (\textquotedblleft{}A deadline; a time limit\textquotedblright{}.) Say it once more.
\end{tcolorbox}
